We present an age-conditioned human motion generation framework for studying aging signatures in the context of assistive robotics.
Our pipeline converts clinical gait recordings into learning-ready motion representations while addressing defects from legacy motion-capture data.
Using 420 clips from 138 subjects across young, mid-age, and elderly groups, we train a Spatial-Temporal Graph Convolutional Network (ST-GCN) to recognize age group from motion clips, achieving 64.6\% classification accuracy.
We then fine-tune a Motion Diffusion Model (MDM) with Low-Rank Adaptation (LoRA) to generate text-conditioned human motion with age-group control.
To evaluate whether generated motions preserve aging-related characteristics, we analyze clinical and kinematic gait metrics on 9,000 generated clips.
Our results show that the model successfully reproduces spatial kinematic constraints in the Elderly condition, but spatiotemporal consistency remains unresolved --- generated stride variability massively overshoots clinical baselines, which we attribute to the capacity ceiling of low-rank, discrete-token adaptation.
These findings highlight dataset size, motion quality, and conditioning strength as key bottlenecks, and motivate larger age-diverse datasets for demographic-aware motion prediction and assistive robot planning.
Code and video available on the \href{https://catr1xliu.github.io/Notebook/Aging-Motion-Project/}{project page}.